% html-boxes.tex -- converter-side definitions of the paper's typed boxes.
%
% WHY: the print versions of these boxes are tcolorboxes, which are drawn with
% TikZ. latexml-oxide converts TikZ faithfully -- by rendering it as SVG. That
% turns 94% of Supplement 1 and 72% of Supplement 2 into fixed-geometry
% pictures: they do not reflow, they do not scale to a narrow screen, and the
% text inside them is not HTML.
%
% So under a converter we skip tcolorbox entirely and define the same
% environments as ordinary theorem-like blocks. LaTeXML and oxide both emit
% these as <div class="ltx_theorem ltx_theorem_NAME"> with a real title, a real
% counter and a real anchor -- semantic HTML that reflows and can be styled.
%
% Input this from inside an \iflatexml branch, then declare the boxes the
% document actually uses:
%
%     \iflatexml
%       % html-boxes.tex -- converter-side definitions of the paper's typed boxes.
%
% WHY: the print versions of these boxes are tcolorboxes, which are drawn with
% TikZ. latexml-oxide converts TikZ faithfully -- by rendering it as SVG. That
% turns 94% of Supplement 1 and 72% of Supplement 2 into fixed-geometry
% pictures: they do not reflow, they do not scale to a narrow screen, and the
% text inside them is not HTML.
%
% So under a converter we skip tcolorbox entirely and define the same
% environments as ordinary theorem-like blocks. LaTeXML and oxide both emit
% these as <div class="ltx_theorem ltx_theorem_NAME"> with a real title, a real
% counter and a real anchor -- semantic HTML that reflows and can be styled.
%
% Input this from inside an \iflatexml branch, then declare the boxes the
% document actually uses:
%
%     \iflatexml
%       % html-boxes.tex -- converter-side definitions of the paper's typed boxes.
%
% WHY: the print versions of these boxes are tcolorboxes, which are drawn with
% TikZ. latexml-oxide converts TikZ faithfully -- by rendering it as SVG. That
% turns 94% of Supplement 1 and 72% of Supplement 2 into fixed-geometry
% pictures: they do not reflow, they do not scale to a narrow screen, and the
% text inside them is not HTML.
%
% So under a converter we skip tcolorbox entirely and define the same
% environments as ordinary theorem-like blocks. LaTeXML and oxide both emit
% these as <div class="ltx_theorem ltx_theorem_NAME"> with a real title, a real
% counter and a real anchor -- semantic HTML that reflows and can be styled.
%
% Input this from inside an \iflatexml branch, then declare the boxes the
% document actually uses:
%
%     \iflatexml
%       % html-boxes.tex -- converter-side definitions of the paper's typed boxes.
%
% WHY: the print versions of these boxes are tcolorboxes, which are drawn with
% TikZ. latexml-oxide converts TikZ faithfully -- by rendering it as SVG. That
% turns 94% of Supplement 1 and 72% of Supplement 2 into fixed-geometry
% pictures: they do not reflow, they do not scale to a narrow screen, and the
% text inside them is not HTML.
%
% So under a converter we skip tcolorbox entirely and define the same
% environments as ordinary theorem-like blocks. LaTeXML and oxide both emit
% these as <div class="ltx_theorem ltx_theorem_NAME"> with a real title, a real
% counter and a real anchor -- semantic HTML that reflows and can be styled.
%
% Input this from inside an \iflatexml branch, then declare the boxes the
% document actually uses:
%
%     \iflatexml
%       \input{html-boxes}
%       \NewHTMLBox{definition}{Definition}{1}   % 1 = takes a {title} argument
%       \NewHTMLBox{theorem}{Theorem}{0}         % 0 = no title argument
%     \else
%       ... the tcolorbox definitions ...
%     \fi
%
% The [label={...}] option each box carries in the sources is honoured, so
% \ref and \autoref into the boxes keep working.

\usepackage{amsthm}
\usepackage{keyval}

% This file is \input, not a package, so `@` is not a letter here: everything
% touching \define@key or \@empty must sit inside \makeatletter.
\makeatletter

% Option keys. Only `label` does anything; the rest are print-layout keys that
% appear in the sources (fold, breakable=false) and must be absorbed silently.
\define@key{htmlbox}{label}{\def\htmlboxlabel{#1}}
\define@key{htmlbox}{fold}[]{}
\define@key{htmlbox}{unfold}[]{}
\define@key{htmlbox}{breakable}[]{}

\newcommand{\htmlboxsetup}[1]{%
  \def\htmlboxlabel{}%
  \def\htmlbox@opts{#1}%
  \ifx\htmlbox@opts\@empty\else\setkeys{htmlbox}{#1}\fi
  \ifx\htmlboxlabel\@empty\else\expandafter\label\expandafter{\htmlboxlabel}\fi}
\makeatother

% \NewHTMLBox{envname}{Heading}{1 if it takes a mandatory {title}, else 0}
\newcommand{\NewHTMLBox}[3]{%
  % amsthm's default `plain' style italicises the body, which is wrong for
  % findings and worked examples that are pages of prose. `definition' style
  % keeps the body upright and only bolds the heading.
  \theoremstyle{definition}%
  \newtheorem{htmlbox#1}{#2}%
  \ifnum#3=1\relax
    \NewDocumentEnvironment{#1}{O{} m}%
      {\begin{htmlbox#1}[##2]\htmlboxsetup{##1}}%
      {\end{htmlbox#1}}%
  \else
    \NewDocumentEnvironment{#1}{O{}}%
      {\begin{htmlbox#1}\htmlboxsetup{##1}}%
      {\end{htmlbox#1}}%
  \fi}

% Present in the print branch; harmless no-ops here so stray uses do not break.
\providecommand{\foldmark}{}
\providecommand{\thetcbcounter}{}
\makeatletter
\@ifundefined{ifpaperboxfold}{\newif\ifpaperboxfold}{}
\makeatother
\endinput

%       \NewHTMLBox{definition}{Definition}{1}   % 1 = takes a {title} argument
%       \NewHTMLBox{theorem}{Theorem}{0}         % 0 = no title argument
%     \else
%       ... the tcolorbox definitions ...
%     \fi
%
% The [label={...}] option each box carries in the sources is honoured, so
% \ref and \autoref into the boxes keep working.

\usepackage{amsthm}
\usepackage{keyval}

% This file is \input, not a package, so `@` is not a letter here: everything
% touching \define@key or \@empty must sit inside \makeatletter.
\makeatletter

% Option keys. Only `label` does anything; the rest are print-layout keys that
% appear in the sources (fold, breakable=false) and must be absorbed silently.
\define@key{htmlbox}{label}{\def\htmlboxlabel{#1}}
\define@key{htmlbox}{fold}[]{}
\define@key{htmlbox}{unfold}[]{}
\define@key{htmlbox}{breakable}[]{}

\newcommand{\htmlboxsetup}[1]{%
  \def\htmlboxlabel{}%
  \def\htmlbox@opts{#1}%
  \ifx\htmlbox@opts\@empty\else\setkeys{htmlbox}{#1}\fi
  \ifx\htmlboxlabel\@empty\else\expandafter\label\expandafter{\htmlboxlabel}\fi}
\makeatother

% \NewHTMLBox{envname}{Heading}{1 if it takes a mandatory {title}, else 0}
\newcommand{\NewHTMLBox}[3]{%
  % amsthm's default `plain' style italicises the body, which is wrong for
  % findings and worked examples that are pages of prose. `definition' style
  % keeps the body upright and only bolds the heading.
  \theoremstyle{definition}%
  \newtheorem{htmlbox#1}{#2}%
  \ifnum#3=1\relax
    \NewDocumentEnvironment{#1}{O{} m}%
      {\begin{htmlbox#1}[##2]\htmlboxsetup{##1}}%
      {\end{htmlbox#1}}%
  \else
    \NewDocumentEnvironment{#1}{O{}}%
      {\begin{htmlbox#1}\htmlboxsetup{##1}}%
      {\end{htmlbox#1}}%
  \fi}

% Present in the print branch; harmless no-ops here so stray uses do not break.
\providecommand{\foldmark}{}
\providecommand{\thetcbcounter}{}
\makeatletter
\@ifundefined{ifpaperboxfold}{\newif\ifpaperboxfold}{}
\makeatother
\endinput

%       \NewHTMLBox{definition}{Definition}{1}   % 1 = takes a {title} argument
%       \NewHTMLBox{theorem}{Theorem}{0}         % 0 = no title argument
%     \else
%       ... the tcolorbox definitions ...
%     \fi
%
% The [label={...}] option each box carries in the sources is honoured, so
% \ref and \autoref into the boxes keep working.

\usepackage{amsthm}
\usepackage{keyval}

% This file is \input, not a package, so `@` is not a letter here: everything
% touching \define@key or \@empty must sit inside \makeatletter.
\makeatletter

% Option keys. Only `label` does anything; the rest are print-layout keys that
% appear in the sources (fold, breakable=false) and must be absorbed silently.
\define@key{htmlbox}{label}{\def\htmlboxlabel{#1}}
\define@key{htmlbox}{fold}[]{}
\define@key{htmlbox}{unfold}[]{}
\define@key{htmlbox}{breakable}[]{}

\newcommand{\htmlboxsetup}[1]{%
  \def\htmlboxlabel{}%
  \def\htmlbox@opts{#1}%
  \ifx\htmlbox@opts\@empty\else\setkeys{htmlbox}{#1}\fi
  \ifx\htmlboxlabel\@empty\else\expandafter\label\expandafter{\htmlboxlabel}\fi}
\makeatother

% \NewHTMLBox{envname}{Heading}{1 if it takes a mandatory {title}, else 0}
\newcommand{\NewHTMLBox}[3]{%
  % amsthm's default `plain' style italicises the body, which is wrong for
  % findings and worked examples that are pages of prose. `definition' style
  % keeps the body upright and only bolds the heading.
  \theoremstyle{definition}%
  \newtheorem{htmlbox#1}{#2}%
  \ifnum#3=1\relax
    \NewDocumentEnvironment{#1}{O{} m}%
      {\begin{htmlbox#1}[##2]\htmlboxsetup{##1}}%
      {\end{htmlbox#1}}%
  \else
    \NewDocumentEnvironment{#1}{O{}}%
      {\begin{htmlbox#1}\htmlboxsetup{##1}}%
      {\end{htmlbox#1}}%
  \fi}

% Present in the print branch; harmless no-ops here so stray uses do not break.
\providecommand{\foldmark}{}
\providecommand{\thetcbcounter}{}
\makeatletter
\@ifundefined{ifpaperboxfold}{\newif\ifpaperboxfold}{}
\makeatother
\endinput

%       \NewHTMLBox{definition}{Definition}{1}   % 1 = takes a {title} argument
%       \NewHTMLBox{theorem}{Theorem}{0}         % 0 = no title argument
%     \else
%       ... the tcolorbox definitions ...
%     \fi
%
% The [label={...}] option each box carries in the sources is honoured, so
% \ref and \autoref into the boxes keep working.

\usepackage{amsthm}
\usepackage{keyval}

% This file is \input, not a package, so `@` is not a letter here: everything
% touching \define@key or \@empty must sit inside \makeatletter.
\makeatletter

% Option keys. Only `label` does anything; the rest are print-layout keys that
% appear in the sources (fold, breakable=false) and must be absorbed silently.
\define@key{htmlbox}{label}{\def\htmlboxlabel{#1}}
\define@key{htmlbox}{fold}[]{}
\define@key{htmlbox}{unfold}[]{}
\define@key{htmlbox}{breakable}[]{}

\newcommand{\htmlboxsetup}[1]{%
  \def\htmlboxlabel{}%
  \def\htmlbox@opts{#1}%
  \ifx\htmlbox@opts\@empty\else\setkeys{htmlbox}{#1}\fi
  \ifx\htmlboxlabel\@empty\else\expandafter\label\expandafter{\htmlboxlabel}\fi}
\makeatother

% \NewHTMLBox{envname}{Heading}{1 if it takes a mandatory {title}, else 0}
\newcommand{\NewHTMLBox}[3]{%
  % amsthm's default `plain' style italicises the body, which is wrong for
  % findings and worked examples that are pages of prose. `definition' style
  % keeps the body upright and only bolds the heading.
  \theoremstyle{definition}%
  \newtheorem{htmlbox#1}{#2}%
  \ifnum#3=1\relax
    \NewDocumentEnvironment{#1}{O{} m}%
      {\begin{htmlbox#1}[##2]\htmlboxsetup{##1}}%
      {\end{htmlbox#1}}%
  \else
    \NewDocumentEnvironment{#1}{O{}}%
      {\begin{htmlbox#1}\htmlboxsetup{##1}}%
      {\end{htmlbox#1}}%
  \fi}

% Present in the print branch; harmless no-ops here so stray uses do not break.
\providecommand{\foldmark}{}
\providecommand{\thetcbcounter}{}
\makeatletter
\@ifundefined{ifpaperboxfold}{\newif\ifpaperboxfold}{}
\makeatother
\endinput
